\section{The tower: centralisers, nesting, and carriers}
\label{sec:tower}

Climbing the tower adds six directions per rung and leaves the local structure of each factor unchanged. The interior four-plane
stays in the first factor; the question is what the added directions do to the internal symmetry.

\begin{theorem}[Internal centraliser]
\label{thm:centraliser}
In $\Cl(3n,3n)$ the bivectors commuting with the interior Lorentz action form
\[
\gint^{(n)}\cong\so{3n-1,3n-3},\qquad \dim\gint^{(n)}=\tfrac12(6n-4)(6n-5).
\]
\end{theorem}

\begin{proof}
By \eqref{eq:bladecomm} a bivector sharing exactly one index with an interior bivector anticommutes with it, and one sharing none
or two commutes. A bivector commuting with all six interior bivectors therefore has either both indices in $E$, which gives
$\lint$ itself and is excluded by commutation with the boosts, or both indices in the Lorentz complement
\[
W_n=\langle e^{(1)}_2,e^{(1)}_3\rangle\oplus\bigoplus_{k=2}^n\langle e^{(k)}_1,\dots,e^{(k)}_6\rangle ,
\]
which has $3n-1$ positive and $3n-3$ negative directions. All bivectors of $W_n$ commute with $\lint$, and they span $\so{W_n}$.
\end{proof}

\begin{corollary}[The photon at one constituent]
\label{cor:photon1}
At $n=1$, $W_1=U$ and $\gint^{(1)}=\R J$. The electromagnetic $\uu1$ is not a symmetry of the interior structure but the
symmetry, by exhaustion. In agreement with \cref{prop:gauge1}, no other charge direction gauges at $n=1$.
\end{corollary}

\begin{proposition}[One frame, and nesting]
\label{prop:frame}
The interior four-plane lies in the first tensor factor, so the six interior bivectors have the same index sets at every rung and
identical structure constants. The tower is one Lorentz frame with an increasing family of modules over it. Moreover
$W_1\subset W_2\subset\cdots$, hence
\[
\gint^{(1)}\subset\gint^{(2)}\subset\gint^{(3)}\subset\cdots,\qquad \dim=1,\ 28,\ 91,\ 190,\ 325,\dots,
\]
with direct limit $\so{\infty,\infty}$. A connection valued at one rung is a connection at every higher rung with no matching
condition. The inclusion runs in the wrong direction to justify a truncation to low rungs, which would be a separate postulate.
\end{proposition}

A gauge algebra has to be compact, and every compact subalgebra of a reductive real Lie algebra is conjugate into a maximal compact
one \cite{Knapp,Helgason}, here $\so{3n-1}\oplus\so{3n-3}$. Since $\su2$ needs a faithful real representation of dimension at least
three and $\su3$ one of dimension at least six \cite{GoodmanWallach}, the rungs can be sorted by what they can carry.

\begin{theorem}[Carrier census]
\label{thm:census}
\[
\begin{array}{c|l|l|r|r|c|c|c}
\toprule
n & \gint^{(n)} & \text{maximal compact} & \dim & \text{rank} & \uu1\ \text{alone} & \su2 & \su3\\
\midrule
1 & \so{2,0} & \so2 & 1 & 1 & \text{yes} & \text{no} & \text{no}\\
2 & \so{5,3} & \so5\oplus\so3 & 13 & 3 & \text{no} & \text{yes} & \text{no}\\
3 & \so{8,6} & \so8\oplus\so6 & 43 & 7 & \text{no} & \text{yes} & \text{yes}\\
4 & \so{11,9} & \so{11}\oplus\so9 & 91 & 9 & \text{no} & \text{yes} & \text{yes}\\
\bottomrule
\end{array}
\]
At $n=2$ the projection of a simple $\su3$ to $\so3$ must vanish and $\so5$ has no proper subalgebra of dimension above six; at
$n=3$ colour enters through $\su3\subset\su4\cong\so6$.
\end{theorem}

\begin{lemma}[Rank obstruction]
\label{lem:rank}
$\rank(\su3\oplus\su2\oplus\uu1)=4$, while the maximal compact subalgebras above have ranks $1,3,7,9$. The full Standard-Model
algebra therefore embeds at no rung $n\le2$. At $n=2$ the maximal compact subalgebra has dimension $13\ge12$, so dimension counting
alone would admit it, and rank is the sharp invariant.
\end{lemma}

\begin{figure}[t]
\centering
\begin{tikzpicture}[scale=1.0]
\foreach \n/\alg/\dimc/\cmp/\cr in {1/{\so{2,0}}/{1}/{\so2}/{\uu1},2/{\so{5,3}}/{28}/{\so5\oplus\so3}/{\uu1,\ \su2},3/{\so{8,6}}/{91}/{\so8\oplus\so6}/{\uu1,\ \su2,\ \su3},4/{\so{11,9}}/{190}/{\so{11}\oplus\so9}/{\uu1,\ \su2,\ \su3}}{
  \pgfmathsetmacro{\y}{1.35*(\n-1)}
  \pgfmathsetmacro{\w}{1.6+0.6*\n}
  \fill[tcol!8] (-\w,\y) rectangle (\w,\y+1.0);
  \draw[tcol,thick] (-\w,\y) rectangle (\w,\y+1.0);
  \node at (0,\y+0.68) {\small $\gint^{(\n)}=\alg$, $\dim\dimc$};
  \node at (0,\y+0.28) {\small compact part $\cmp$};
  \node[anchor=west] at (4.3,\y+0.5) {\small $n=\n$: carries $\cr$};
}
\foreach \n in {1,2,3}{\pgfmathsetmacro{\y}{1.35*\n-0.35}\draw[->,thick,gray] (0,\y-0.0) -- (0,\y+0.35);}
\node[anchor=north] at (0,-0.15) {\small nested: $\gint^{(1)}\subset\gint^{(2)}\subset\gint^{(3)}\subset\cdots$};
\end{tikzpicture}
\caption{The tower of internal centralisers (\cref{thm:centraliser,thm:census}). Each rung contains the previous one; the compact part
decides which gauge factors can be carried. The electromagnetic $\uu1$ is the whole centraliser only at the first rung, the one rung
with no room for a nonabelian factor.}
\label{fig:tower}
\end{figure}

\begin{corollary}[The tower's trade-off]
\label{cor:tradeoff}
The $\uu1$ is the whole internal centraliser only at $n=1$, the one rung with no room for either nonabelian carrier. Inside
$\gint^{(n)}$ the centraliser of the embedded one-rung charge $e^{(1)}_2e^{(1)}_3$ has dimensions $1,16,67,154,277$ for $n=1,\dots,5$,
since by \eqref{eq:bladecomm} exactly $12(n-1)$ internal bivectors fail to commute with it. From $n=2$ onwards the photon is one
direction in a space that grows without bound; its uniqueness is lost to no dynamical process, only to the arrival of room.
\end{corollary}

At $n\ge2$ the embedded one-rung charge $e^{(1)}_2e^{(1)}_3$ must not be confused with the complex structure $J=\wI$ of
\cref{thm:i}, which commutes with every internal bivector and has all of $\gint^{(n)}$ as centraliser; they agree only at $n=1$.

\begin{remark}[The gaugeable fraction does not improve]
Of the $\binom{6n}2=3n(6n-1)$ bivectors of $\Cl(3n,3n)$, $2\binom{3n}2=3n(3n-1)$ pair directions of the same sign: the fractions
$6/15$, $30/66$, $72/153$ rise only to $1/2$. Climbing adds boosts at essentially the rate at which it adds rotations. The census
permits the chain $\uu1\subset\uu1\oplus\su2\subset\uu1\oplus\su2\oplus\su3$ at $n=1,2,3$; it permits it, and nothing selects those
subalgebras among others of the same rank.
\end{remark}
